\begin{itemize}
    \item $A_0$ - The Sensible Object in Actuality (\gk{ἐνεργείᾳ} \gk{αἰσθητόν} / \textit{energeiai aisthēton})
    \begin{itemize}
        \item The sensible quality is actually exercising its power---the colour is visible, the sound audible---such that the object functions as the \textit{first} actuality in the chain, the unmoved originator that sets the entire perceptual event into motion without itself undergoing any alteration through the act of being perceived (\textit{DA} II.5, 417a-b). Indeed, the sensible object in actuality occupies a position structurally analogous to that of the unmoved mover in Aristotle's broader metaphysics: it initiates a causal sequence by exercising a power that is already fully realized, and it does so without suffering any reciprocal change from the motion it occasions.
        \item $M_0 \rightarrow M_1$: Perceptual Motion (\gk{αἴσθησις} / \textit{aisthēsis})
        \begin{itemize}
            \item The fulfillment of the sense organ's potential \textit{qua} capable of receiving the respective sensible form---the transition, that is, from the organ's dormant first actuality (its standing capacity to be affected) to the exercise of that capacity in the presence of the sensible object.
            \item Kinetic Roles:
            \begin{itemize}
                \item Unmoved Originator: $A_0$---the sensible object in actuality.
                \item Moved Mover: the sense faculty, which receives the form without the matter and is thereby altered (\textit{DA} II.12, 424a17-24).
                \item That Which is Moved: the sense organ and the perceiver as material substrate.
                \item Medium: air, water, or flesh---that through which the sensible form is transmitted from object to organ.
                \item Time ($T_0 \rightarrow T_1$): the temporal interval over which the medium transmits and the organ receives---nearly simultaneous for sight, yet measurably extended for sound, whose propagation through the medium Aristotle treats as a genuinely successive process.
            \end{itemize}
        \item This motion is named \textit{perception} because its terminus is $A_1$. The affective dimension, it must be emphasized, is constitutive rather than secondary: sensible form and affective valence arrive together, as a single indivisible event, in the moment at which the world impresses itself upon the soul.
        \item Destruction condition (426a27-b8): if the sensible exceeds the ratio of the organ, \hl{the motion acts on the organ \textit{as matter} rather than \textit{as sense}.} The ratio is destroyed, and $A_1$ is accordingly never reached. The excess names no completed actuality and therefore constitutes no genuine perceptual motion; it is damage, not perception. \hl{What Aristotle describes here is a failed actualization of the sensory power, not the absence of motion tout court. There is real change---genuine damage---in the organ, but it is not the right sort of change: it is no genuine perceptual motion in the sense that there is no motion qua aisthēsis, no transition from potentiality to the determinate actuality of perceiving.}
        \end{itemize}
    \end{itemize}
\end{itemize}

\begin{itemize}
    \item $A_1$ - Completed Perception: Aisthēsis and Aisthēma (\gk{ἐνέργεια αἰσθήσεως} / \textit{energeia aistheseos})
    \begin{itemize}
        \item The sense has received the form without the matter, and the ratio (\gk{λόγος} / \textit{lógos}) of the organ has been actualized within its proper range. This constitutes a \textit{second actuality}: the organ possessed the capacity as a first actuality---the sleeping geometer who holds geometrical knowledge without exercising it---and has now passed into the exercise of that capacity, as the geometer who actively proves a theorem. \hl{The actualization of the sensible object and that of the sensitive power are, Aristotle insists, one and the same activity, though they differ in their being---a single event described under two logoi.}
        \item The \textit{aisthēma} is not merely a formal imprint stamped upon a passive substrate; it is, rather, the perceptual act completed, wherein both the sensible form and the soul's affective response (\gk{πάθος} / \textit{pathos}) are present as a unity that admits no internal division. Aristotle observes that ``to perceive then is like bare asserting or thinking; but when the object is pleasant or painful, the soul makes a sort of affirmation or negation, and pursues or avoids the object'' (\textit{DA} III.7, 431a8--10). Thus the perceptual faculty and the appetitive faculty are not two discrete capacities subsequently conjoined but one activity distinguished only in account: ``Both avoidance and appetite when actual are identical with this: the faculty of appetite and avoidance are not different, either from one another or from the faculty of sense-perception; but their being is different'' (431a13--14). The connection, accordingly, is analytic rather than synthetic: wherever there is perception, there is necessarily pleasure and pain, and wherever these are found, there is necessarily appetite (\textit{DA} II.3, 414b4--6). The world has disclosed itself to the perceiver in a determinate way---as both \textit{what it is} and \textit{how it matters}---and this twofold disclosure is accomplished in a single act rather than through the sequential conjunction of a cognitive and an evaluative moment.
\inlinenote{In III.3, Aristotle uses aisthemata (in the plural) as what remain in the organs and can be re-activated, described as impression or picture left by the movement of perception. The dispositional element is textually backed by the analogy with knowledge vs. its exercise (III.2-3), where an acquired hexis can be re-activated.}
        \item $A_1$ names $M_0 \rightarrow M_1$. It also serves as the unmoved originator of $M_1 \rightarrow M_2$---a dual role that is structurally essential, for what completes the first motion simultaneously occasions the second.
\inlinenote{II.5-6: the perceptual motion runs from object, through medium, to organ; the actuality of sense is the end-state. III.3: the aisthemata left behind become the basis for imagination; phantasia is defined as kinesis resulting from ``actual exercise of a power of sense,'' and that movement cannot exist apart from sensation.}
        \item $M_1 \rightarrow M_2$: \textit{Phantastic} Motion: From Residual Trace to \textit{Phantasma}
        \begin{itemize}
            \item After the sensible object withdraws, residual motion persists in the sensory apparatus (\textit{DA} III.3, 428b10--16; \textit{De Insomniis} 459a--461b). This persistence is not a second, discrete event but a single continuous process: the \textit{kinēsis} produced by perception reverberates within the organ and gradually settles into a determinate \textit{phantasma}. The residual motion already carries both dimensions of the original perceptual act---the formal imprint (\textit{aisthēma}) and the affective charge (\textit{pathos})---precisely because it inherits the constitutive unity of the act from which it derives. There is, accordingly, no moment of fusion between a cognitive trace and an affective overlay; there was never any separation to begin with.
            \item Kinetic Roles:
            \begin{itemize}
                \item Unmoved Originator: $A_1$---the completed perception, with its resonant \textit{aisthēma} and \textit{pathos}, which now functions as the source of a new motion without itself being further altered by it.
                \item Moved Mover: \textit{phantasia} (\gk{φαντασία}) as the retentive faculty, which sustains and transforms the residual impression and, in so doing, is itself affected---carrying forward the affection (\gk{πάθος} / \textit{pathos}) that the original perceptual act deposited. \hl{A crucial distinction must be maintained: \textit{phantasia} (the faculty or activity) and \textit{phantasma} (the image) are not identical. Imagination is not the image; it is the \textit{kinēsis} arising from sense ``in virtue of which'' images come into being. Phantasia sustains, activates, and transforms the residual impression into a determinate image---it is the ongoing motion, not the settled product.}
                \item That Which is Moved: the psycho-somatic state of the animal, qualitatively changed by the retained impression. \hl{The description is psycho-somatic or, more precisely, psychophysical: Aristotle consistently treats sensation, imagination, memory, and the affections as ``enmattered accounts'' (\gk{ἔνυλοι λόγοι})---as operations of ``soul and body together'' that resist any clean separation into mental and bodily components.}
                \item Time ($T_1 \rightarrow T_2$): the interval between the withdrawal of the sensible object and the settling of the residual trace into a determinate \textit{phantasma}---brief for afterimages, yet considerably extended for powerful impressions that settle slowly as the motion reverberates through the sensory apparatus.
            \end{itemize}
        \item Named by its terminus: this is \textit{phantastic motion} because it terminates in $A_2$, the crystallization of the residual trace into a stable, recallable image.
        \end{itemize}
    \end{itemize}
\end{itemize}

\begin{itemize}
      \item $A_2$ --- The \textit{Phantasma} Proper (\gk{φάντασμα})
      \begin{itemize}
          \item The determinate image has emerged: the residual motion has settled
          into a dual-aspect unity of formal imprint (\textit{aisthēma}) and
          affective charge (\textit{pathos}), a unity in which the sensible form of
          the object and the soul's dispositional orientation toward it achieve
          determinate expression. The \textit{phantasma}, thus constituted, possesses
          three characteristics essential to its role in the actualization chain:
          \begin{itemize}
              \item It is stable and recallable---available, that is, as an object
              for further cognitive operations that may occur at considerable
              temporal remove from the original perceptual encounter.
              \item It is formally similar to the original perception yet operates
              in the absence of the sensible object (\textit{DA} III.3, 428b10--16:
              ``imagination is held to be a movement \ldots\ necessarily similar in
              character to the sensation itself'')---a similarity that is the source
              both of its cognitive utility and of its susceptibility to error.
              \item It serves as the indispensable hinge between perception and
              thought, for Aristotle insists that ``the soul never thinks without an
              image'' (\textit{DA} III.7, 431b2)---a claim that binds the highest
              operations of \textit{nous} to the phantasmatic residue of sensory
              experience.
          \end{itemize}
          \item $A_2$ names $M_1 \rightarrow M_2$, and it simultaneously serves as
          the unmoved originator of $M_2 \rightarrow M_3$. This is, accordingly, the
          pivotal actuality in the entire chain: everything prior converges upon it;
          everything subsequent operates upon it as its indispensable material.

          \item $M_2 \rightarrow M_3$: \textbf{Noetic / Doxastic / Mnemonic / Anticipatory
          Motion: Cognitive Engagement with the \textit{Phantasma}}
          \begin{itemize}
              \item The \textit{phantasma} is taken up by \textit{nous}, the doxastic
              faculty, memory, or deliberative \textit{phantasia}---four parallel
              modes of cognitive engagement that are distinguished not by their
              position in the causal chain but by their intentional direction,
              temporal orientation, and relationship to truth:
\inlinenote{Aristotle distinguishes sensitive and deliberative phantasia (433b29--434a10) but never explicitly says they are parallel orientations of the same apparatus. The unification is implicit in the texts. Likewise, he never says that temporal awareness is what makes possible deliberative phantasia; the capacity to calculate and the capacity to perceive time seem to be co-implicated. Bowin argues that deliberative phantasia actively manipulates images to produce determinate temporal representations---treating it as a function of deliberative phantasia, not its condition.}
              \begin{enumerate}
                  \item \textbf{Intellection}: the \textit{phantasma} addressed
                  \textit{as bearing a universal}---\textit{nous} abstracts the
                  intelligible form from the sensory image, apprehending what is
                  common across particulars. Temporal orientation: \textit{atemporal}.
                  The universal, once grasped, is not indexed to past, present, or
                  future; it holds across all instances and is accordingly not
                  situated within the temporal continuum that conditions the other
                  modes.

                  \item \textbf{Doxa} (\gk{δόξα}): the \textit{phantasma} addressed
                  \textit{as warranting propositional assertion}---the soul commits
                  itself, through an act that is at once cognitive and involuntary, to
                  a truth-claim about what is the case. Doxa constitutes the
                  \textit{logos}-dependent, assertoric operation upon the
                  \textit{phantasma} that phantasia itself is not, for Aristotle
                  insists that ``imagination does not involve opinion based on
                  inference, though opinion involves imagination'' (\textit{DA} III.11,
                  434a10--11). It requires three conditions that phantasia
                  characteristically lacks: belief (\gk{πίστις} / \textit{pistis}),
                  conviction (\gk{πεπεῖσθαι} / \textit{pepeisthai}), and discourse of
                  reason (\gk{λόγος}). Aristotle states this requirement explicitly:
                  ``opinion involves belief, belief involves being convinced, and being
                  convinced involves reason; but some of the brutes have imagination,
                  without discourse of reason'' (\textit{DA} III.3, 428a19--24).
                  Accordingly, doxa is found only in rational animals---those beings
                  that possess the calculative faculty through which
                  \textit{logos} operates.

                  Doxa differs from the other modes of cognitive engagement in two
                  decisive respects:

                  \begin{itemize}
                      \item \textit{Involuntariness}: Aristotle observes that
                      ``imagining lies within our power whenever we wish \ldots\ but
                      in forming opinions we are not free: we must either be holding
                      what is true or what is false'' (\textit{DA} III.3, 427b17--21).
                      Phantasia is \gk{ἐφ' ἡμῖν}; doxa is \gk{οὐκ ἐφ' ἡμῖν}. One
                      can summon any image at will---calling up a picture, practising
                      the mnemonic art---but one cannot choose what one takes to be
                      the case. Doxa arises from the way things appear under the sway
                      of phantasia, from the taking-as function through which the
                      phantasma discloses its object as something determinate; yet once
                      it arises, the soul is bound by its content, constrained by the
                      truth-relation that doxa inherently bears. The involuntariness
                      is thus grounded externally: doxa is answerable to how things
                      are, not to how we might wish them to appear.
                      \item \textit{Independence from the phantasma's own testimony}:
                      doxa and phantasia can deliver contradictory content
                      simultaneously, a phenomenon that demonstrates their ontological
                      distinctness with particular clarity. Aristotle provides the
                      canonical illustration: ``We imagine the sun to be a foot in
                      diameter though we are convinced (\gk{πεπεῖσθαι}) that it is
                      larger than the inhabited earth'' (\textit{DA} III.3, 428b2--4).
                      The \textit{phantasma} is not corrected by the doxa---it
                      persists unchanged, presenting its testimony regardless of what
                      reason concludes---but neither is doxa compelled by the
                      \textit{phantasma}'s apparent testimony. Doxa is thus a distinct
                      assertoric act that takes the \textit{phantasma} as its material
                      yet can override or diverge from its content. This is precisely
                      why Aristotle concludes that ``a true opinion becomes false only
                      when the fact alters without being noticed; imagination is
                      therefore neither any one of the states enumerated, nor
                      compounded out of them'' (428b4--10).
                  \end{itemize}

                  The formation of doxa is, furthermore, subject to a
                  \textit{gatekeeping condition}: whether the corrective rational
                  faculty intercepts the phantasma before doxastic commitment occurs.
                  When the controlling faculty operates normally, the phantasma may be
                  corrected---as in the sun case, where we see ``one foot'' yet do not
                  believe it. When, however, the corrective faculty is disabled---by
                  sleep, illness, passion, or perceptual distance---the phantasma
                  passes unchallenged into doxa, and the soul is thereby committed to
                  the appearance as though it were the fact. Aristotle describes this
                  vulnerability in the \textit{De Insomniis}: ``when the controlling
                  sense is not disabled \ldots\ we are not deceived; but when it is
                  disabled, as happens in sleep and illness, then what appears is taken
                  to be what is the case'' (\textit{De Insomniis} 460b3--16; cf.\
                  \textit{DA} III.3, 428b25--429a2, where Aristotle notes that
                  phantasia is ``especially'' erroneous ``when the object of perception
                  is far off''). The speed of this transition, moreover, can be
                  near-instantaneous: ``when we take something to be fearful
                  (\gk{δοξάσωμεν}), emotion is immediately produced
                  (\gk{εὐθὺς πάσχομεν})'' (\textit{DA} III.3, 427b21--24).

                  Temporal orientation: primarily \textit{present-situational}---doxa
                  asserts ``this is so'' or ``this is not so'' with respect to what the
                  phantasma discloses---though doxa may equally concern past or future
                  states of affairs. Its temporal character is accordingly assertoric
                  rather than directional: doxa \textit{commits} the soul to a
                  truth-claim without privileging any particular temporal horizon.

\inlinenote{The relationship between doxa and the practical syllogism is structurally significant. In \textit{MA} 701a7--25, the practical syllogism proceeds from a universal premise (which may be a settled doxa or epistēmē: ``all sweet things are to be tasted'') conjoined with a particular perception or phantasma (``this is sweet'') to issue in action. Doxa can thus function both as an \textit{output} of the current perceptual event (the particular judgment ``this is dangerous'') and as a pre-existing \textit{hexis} (a settled universal belief) that informs how subsequent phantasmata are processed. In the latter case, the universal doxa operates as an additional unmoved originator alongside $A_2$: the repertoire of settled beliefs shapes how the phantasma is taken up in $M_2 \rightarrow M_3$.}

                  \item \textbf{Memory}: the \textit{phantasma} addressed \textit{as
                  a likeness of something past}---the past is marked as past and
                  felt as past, recognized not merely as a content that once occurred
                  but as belonging to a time that has elapsed (\textit{On Memory}
                  449b24--25). Temporal orientation: \textit{retrospective}. Memory is
                  constitutively temporal in a way that intellection is not; its object
                  \textit{is} the past \textit{qua} past, and the temporal distance
                  between the original perception and the present act of recall is
                  internal to what memory discloses rather than merely external to it.

                  \item \textbf{Deliberative / Anticipatory \textit{Phantasia}}
                  (\gk{βουλευτικὴ φαντασία}): \textit{several} \textit{phantasmata}
                  synthesized into a unity and addressed \textit{as composable into
                  unrealized possibilities}, measured against a single standard of
                  evaluation. Aristotle specifies this operation with considerable
                  precision: ``whether this or that shall be enacted is already a task
                  requiring calculation; and there must be a single standard to measure
                  by, for that is pursued which is greater. It follows that what acts
                  in this way must be able to make a unity out of several images''
                  (\textit{DA} III.11, 434a5--10). Temporal orientation:
                  \textit{prospective}. This mode operates through two interlocking
                  sub-functions:
                  \begin{itemize}
                      \item \textit{Projective-synthetic}: the \textit{phantastic}
                      faculty recombines elements drawn from the available repertoire
                      of \textit{phantasmata} into images of states not yet
                      actual---possible futures, anticipated outcomes, unrealized
                      configurations. This is the productive capacity that goes beyond
                      mere retention: \textit{phantasia} here constructs rather than
                      merely preserves, assembling from the materials of past
                      experience images of what might yet come to pass.
                      \item \textit{Calculative-practical}: the projected
                      possibilities are measured against one another by a single
                      standard (\gk{ἑνί τινι δεῖ μετρεῖν}) to determine which
                      course of action shall be enacted. It is here that
                      \textit{phantasia} converges with practical thought
                      (\gk{διάνοια πρακτική}): the projected image is evaluated not
                      merely as possible but as \textit{to-be-pursued} or
                      \textit{to-be-avoided}, and the deliberative process terminates
                      in a determination of the greater good.
                  \end{itemize}
                  This dual structure is precisely what distinguishes rational action
                  from appetitive reaction. Non-rational animals act from sensitive
                  \textit{phantasia} operating upon a single present
                  \textit{phantasma}---``I want to drink, says appetite; this is
                  drink, says sense or imagination: straightaway I drink''
                  (\textit{MA} 701a32--33). Rational animals, by contrast, because
                  they possess \gk{βουλευτικὴ φαντασία}, can hold back from what is
                  immediately present by projecting what is future---and this capacity
                  requires, indeed presupposes, a \textit{sense of time}: ``while
                  thought bids us hold back because of what is future, desire is
                  influenced by what is just at hand: a pleasant object which is just
                  at hand presents itself as both pleasant and good, without condition
                  in either case, because of want of foresight into what is farther
                  away in time'' (\textit{DA} III.10, 433b5--10). The sense of time is
                  thus co-implicated with deliberative \textit{phantasia}: the capacity
                  to calculate and the capacity to perceive time cannot, on Aristotle's
                  account, be separated without destroying both.

\inlinenote{This is why thought can override desire when considering future implications. Pursuit of the good still drives all movement, but in a calculative, prospective fashion.}

              \end{enumerate}

              \item These modes may operate simultaneously or in succession, for they
              are not stages arranged in a fixed sequence but orientations of the same
              cognitive apparatus toward different aspects of its object. Memory is not
              a separate motion-stage but a specific orientation of the phantasmatic
              capacity toward the temporal past. Likewise, anticipatory
              \textit{phantasia} is not a separate faculty but the
              \textit{prospective} orientation of the same \textit{phantastic}
              capacity that, turned backward, yields memory and, turned inward toward
              the intelligible content of the image, furnishes \textit{nous} with the
              material it requires for abstraction. Doxa, for its part, constitutes the
              \textit{assertoric} engagement with the \textit{phantasma}---the soul's
              involuntary commitment to the appearance as true or false---which may
              accompany or inflect any of the other modes. One may remember \textit{and}
              simultaneously hold a belief about what is remembered; one may deliberate
              \textit{and} form opinions about the projected outcomes over which one
              deliberates. Doxa is thus not a separate temporal orientation but a
              truth-committing layer that can supervene upon any of the three
              orientations---atemporal, retrospective, or prospective---adding to each
              the dimension of involuntary assertoric commitment.

              \item Kinetic Roles:
              \begin{itemize}
                  \item Unmoved Originator: $A_2$---the \textit{phantasma} (or, in
                  the case of deliberative \textit{phantasia}, the available
                  \textit{repertoire} of \textit{phantasmata}) functioning as the
                  object of thought (\textit{DA} III.7, 431b2). For doxa specifically,
                  settled universal beliefs (\textit{doxai} held as \textit{hexeis})
                  may also function as additional unmoved originators, shaping how the
                  \textit{phantasma} is taken up and processed (cf.\ the universal
                  premise in the practical syllogism, \textit{MA} 701a7--25).
                  \item Moved Mover: \textit{nous} (\gk{νοῦς}) for intellection; the
                  doxastic faculty for opinion-formation, operating through
                  \gk{λόγος}, \gk{πίστις}, and \gk{πεπεῖσθαι} (\textit{DA} III.3,
                  428a19--24); the mnemonic faculty for memory;
                  \gk{βουλευτικὴ φαντασία} for deliberation and anticipation---each
                  activated by and operating upon the \textit{phantasma} as its
                  proximate object.
                  \item That Which is Moved: the knower, believer, rememberer,
                  deliberator, or anticipator---the rational animal, cognitively
                  transformed by the engagement.
                  \item Time ($T_2 \rightarrow T_3$): variable according to mode.
                  Intellection of universals can be instantaneous, for the universal
                  is atemporal in character and is grasped in a single act rather than
                  extended over duration. Doxa-formation varies in its temporal
                  extension: simple perceptual doxa (``this is white,'' ``this is
                  dangerous'') may be near-instantaneous, its speed indicated by
                  Aristotle's \gk{εὐθύς}---``immediately'' (\textit{DA} III.3,
                  427b23)---whereas inferential doxa, which involves combining
                  premises through \textit{logos}, extends over the duration of the
                  reasoning process. Memory is constitutively temporal: its object
                  \textit{is} the past \textit{qua} past, and the temporal interval
                  between the original perception and the present act of recall is
                  internal to what memory discloses. Deliberative and anticipatory
                  \textit{phantasia} is constitutively \textit{prospective}: it
                  extends over the weighing of alternatives projected into the
                  future, and its temporal horizon is bounded by the ``sense of
                  time'' (\textit{DA} III.10, 433b7) that enables foresight into
                  what is farther away. Time here is thus not merely the measure of
                  the motion but the medium within which the four modes differentiate
                  themselves by intentional direction: beyond duration (intellection),
                  assertoric commitment (doxa), backward (memory), and forward
                  (anticipation).
              \end{itemize}

              \item Named by its terminus: this is
              \textit{thinking / opining / remembering / deliberating / anticipating}
              because it terminates in $A_3$.
          \end{itemize}
      \end{itemize}
  \end{itemize}

  \begin{itemize}
      \item $A_3$ --- Completed Noetic / Doxastic / Mnemonic / Anticipatory Actuality
      \begin{itemize}
          \item The judgment is formed. The universal has been grasped, the opinion
          is held, the past has been recalled and recognized as past, the
          deliberation has reached its conclusion, or the anticipated possibility
          has been constructed, evaluated, and either adopted or rejected. The
          \textit{phantasma}---or, in the case of deliberative \textit{phantasia},
          the synthesized unity of several \textit{phantasmata}---has been fully
          processed, refined in both its formal content and its felt significance,
          such that the cognitive engagement initiated at $M_2$ has achieved its
          determinate completion.

          \item In the case of doxa, $A_3$ is a determinate belief: the soul now
          holds that something \textit{is} or \textit{is not} the case, and this
          holding is not a provisional entertainment but an involuntary commitment
          from which the subject cannot freely withdraw. This completed opinion
          differs from the completed acts of the other modes in two structurally
          decisive respects: its involuntary bindingness---once formed, the subject
          is not free to discard it at will (\textit{DA} III.3, 427b17--21)---and
          its distinctive affective consequence. Where the \textit{phantasma} at
          $A_2$ carries latent affective valence (\textit{pathos} as dispositional
          orientation toward pursuit or avoidance), doxa \textit{activates} that
          valence into a condition capable of producing higher-order emotion.
          Aristotle draws the contrast with characteristic precision: ``when we merely
          imagine something fearful or encouraging we remain unaffected
          (\gk{ἀπαθῶς ἔχομεν}), but when we take something to be fearful
          (\gk{δοξάσωμεν}), emotion is immediately produced
          (\gk{εὐθὺς πάσχομεν})'' (\textit{DA} III.3, 427b21--24). The
          \textit{phantasma}'s latent \textit{pathos} is accordingly necessary but
          not sufficient for the mobilization of the soul; it is doxa's
          truth-commitment---the soul's involuntary taking-as-true---that converts
          dispositional orientation into the mobilized affective state from which
          higher-order emotions (fear, anger, desire, and the rest) can arise. The
          analysis of emotion as a distinct moment in the chain, with its own
          internal structure and its own position in the actualization sequence, is
          reserved for subsequent treatment.

\inlinenote{The distinction between \textit{pathos} (affective valence present from $A_1$ onward) and \textit{emotion} (the higher-order affective states---fear, anger, pity, etc.---that arise only after doxastic commitment) is structurally load-bearing. Pathos is the raw good/bad, pleasant/painful orientation impressed with the formal trace at the moment of perception and inherited by the phantasma. Emotion is the determinate mobilization of the soul that follows from \textit{taking-as-true}: the soul is not merely disposed toward avoidance (pathos) but is \textit{afraid} (emotion). The passage at 427b21--24 names precisely this transition: \gk{ἀπαθῶς ἔχομεν} under bare phantasia; \gk{εὐθὺς πάσχομεν} under doxa. The emotion moment will be developed in a subsequent revision of the chain.}

          \item What emerges from $A_3$ is the apprehension of a realizable good (or
          evil to be avoided): the object of desire that initiates the final motion
          toward action. Judgment and affective evaluation have converged---either
          explicitly through deliberation or compressed, through long habituation,
          into kinesthetic schemata of practiced expertise. In the case of
          deliberative \textit{phantasia}, the ``realizable good'' is not merely
          recognized in what is present but \textit{projected} into an anticipated
          future and selected from among competing possibilities by a single standard
          of measure. In the case of doxa, the realizable good (or the evil to be
          avoided) is disclosed through the soul's involuntary commitment to the
          appearance as true: the object is not merely imagined as threatening but
          \textit{taken to be} threatening, and it is precisely this taking-as-true
          that gives it the motivational force requisite for initiating movement.
          Speculative doxa, however, need not issue in action at all, for Aristotle
          distinguishes practical thought, ``which calculates means to an end,'' from
          speculative thought, which ``moves nothing'' (\textit{DA} III.10,
          433a13--15). Where doxa is speculative---where the soul holds a belief about
          what is the case without that belief disclosing a realizable good or
          evil---$A_3$ does not serve as the unmoved originator of
          $M_3 \rightarrow M_4$, and the chain accordingly terminates at $A_3$ for
          that mode.

          \item $A_3$ names $M_2 \rightarrow M_3$. It serves as the unmoved
          originator of $M_3 \rightarrow M_4$---but only mediately, through the
          realizable good it discloses. The good that moves appetite may be
          immediately present (sensitive \textit{phantasia}: the drink before me),
          anticipatorily projected (deliberative \textit{phantasia}: the course of
          action whose outcome I foresee), or involuntarily taken to be the case
          (doxa: the thing believed dangerous). In each case, it is $A_3$ as
          completed actuality---not the raw \textit{phantasma} at $A_2$---that
          functions as the unmoved originator of desire.
\inlinenote{433b10--18 describes the object of desire as that which moves without being moved, and desire as that which moves and is moved.}
          \item $M_3 \rightarrow M_4$: \textbf{Orektikon Motion: Desire Converging
          in Action}
          \begin{itemize}
              \item Judgment and desire converge in overt action, drawing upon the
              full kinetic-affective history of the \textit{phantasma} as it has been
              elaborated through the preceding stages of the chain. Where deliberative
              \textit{phantasia} has operated, the desire is informed by foresight:
              the animal acts not merely from what is at hand but toward what has been
              projected and evaluated as the greater good. Where only sensitive
              \textit{phantasia} is at work, appetite responds to the immediately
              present \textit{phantasma} without foresight---``because of want of
              foresight into what is farther away in time'' (\textit{DA} III.10,
              433b9--10). Where doxa has operated, the desire is informed by the
              soul's truth-commitment: the animal acts not merely from how things
              appear but from how they are \textit{taken to be}. The affective force
              of doxa-mediated desire is, Aristotle indicates, both immediate
              (\gk{εὐθύς}, \textit{DA} III.3, 427b23) and involuntary
              (\gk{οὐκ ἐφ' ἡμῖν}, 427b20)---the soul that believes something
              dangerous is already mobilized toward avoidance before deliberation can
              intervene, precisely because doxa, unlike phantasia, is not subject to
              the will.
              \item Kinetic Roles:
              \begin{itemize}
                  \item Unmoved Originator: the realizable good as apprehended in
                  $A_3$ (\textit{DA} III.10, 433b10--18).
                  \item Moved Mover: appetite or desire (\gk{ὄρεξις} /
                  \textit{orexis})---``that which at once moves and is moved''
                  (433b17--18); ``appetite in the sense of actual appetite
                  \textit{is} a kind of movement'' (433b18). Desire is thus the
                  intermediate kinetic element that translates the apprehended good
                  into bodily motion.
                  \item That Which is Moved: the animal via the bodily
                  instrument---the ball-and-socket joint where ``pushing and
                  pulling'' execute the action (433b21--25).
                  \item Time ($T_3 \rightarrow T_4$): the interval between the
                  formation of desire and the completion of the motor act---an
                  interval whose duration varies with the complexity of the action
                  and the resistance of the material conditions.
              \end{itemize}
              \item Named by its terminus: this is \textit{appetitive action}
              because it terminates in $A_4$.

\inlinenote{the source of all animal movement. MA CITE}

          \end{itemize}
      \end{itemize}
  \end{itemize}

  \begin{itemize}
      \item $A_4$ --- Completed Action (\gk{πρᾶξις})
      \begin{itemize}
          \item The perceptual-\textit{phantasmatic}-motor loop is closed, and the
          ensouled body has engaged the moving world. The completed action draws upon
          the full kinetic-affective history of the \textit{phantasma}---from the
          original perceptual encounter through its retention as residual trace,
          its cognitive reactivation in the modes of thought, belief, memory, and
          anticipation, and its appetitive appropriation as a realizable good that
          moves the animal to act.
          \item $A_4$ names $M_3 \rightarrow M_4$. The completed action, however,
          does not merely close the chain; it changes the world, producing new
          sensible conditions that can serve as $A_0$ for a new perceptual cycle.
          The chain is thus not linear but recursive: each completed action
          reconstitutes the perceptual environment and thereby occasions a new
          sequence of actualizations, a new passage from the sensible object through
          phantasma, cognition, desire, and movement.
      \end{itemize}
  \end{itemize}
